\documentclass[11pt]{article}

\usepackage[margin=1in]{geometry}
\usepackage{amsmath,amssymb,mathtools}
\usepackage{bm}

\newcommand{\AdS}{\mathrm{AdS}}
\newcommand{\dd}{\mathrm{d}}
\newcommand{\ii}{\mathrm{i}}
\newcommand{\pd}{\partial}
\newcommand{\cF}{\mathcal{F}}
\newcommand{\cD}{\mathcal{D}}

\title{Finite-Radius Scalar and Vector Systems in Global \(\AdS_3\)}
\author{}
\date{}

\begin{document}
\maketitle

\section{Setup}

Use global \(\AdS_3\)
\begin{align}
\dd s^2=-(1+r^2)\dd t^2+\frac{\dd r^2}{1+r^2}+r^2\dd\phi^2,
\qquad 0\le r\le r_0 .
\end{align}
The finite wall is placed at \(r=r_0\).  For the Maxwell system we impose
\begin{align}
A^t\big|_{r_0}=0,\qquad A^\phi\big|_{r_0}=0 ,
\end{align}
and work in temporal gauge \(A^t=0\).  Regularity at \(r=0\) fixes the
regular radial branch.

\section{Scalar Field}

For a free scalar
\begin{align}
S=-\frac12\int\dd^3x\sqrt{-g}\left(\nabla_\mu\phi\nabla^\mu\phi
+M^2\phi^2\right),
\end{align}
write the separated modes as
\begin{align}
\phi(t,r,\phi)=e^{-\ii\omega t+\ii m\phi}R_{\omega m}(r).
\end{align}
The radial equation is
\begin{align}
(1+r^2)R_{\omega m}''+\frac{1+3r^2}{r}R_{\omega m}'
+\left(\frac{\omega^2}{1+r^2}-\frac{m^2}{r^2}-M^2\right)R_{\omega m}=0.
\end{align}
When \(M^2=\Delta(\Delta-2)\), the branch regular at the origin is
\begin{align}
R_{\omega m}^{(\Delta)}(r)=
r^{|m|}(1+r^2)^{-(|m|+\Delta)/2}
{}_2F_1\!\left(
\frac{|m|+\Delta-\omega}{2},
\frac{|m|+\Delta+\omega}{2};
1+|m|;
\frac{r^2}{1+r^2}
\right).
\end{align}

\subsection{Dirichlet Wall and Standard Quantization}

For standard quantization take
\begin{align}
\Delta=\Delta_+=1+\sqrt{1+M^2}>1 .
\end{align}
The finite-radius Dirichlet condition is
\begin{align}
R_{\omega m}^{(\Delta)}(r_0)=0,
\end{align}
or equivalently
\begin{align}
{}_2F_1\!\left(
\frac{|m|+\Delta-\omega}{2},
\frac{|m|+\Delta+\omega}{2};
1+|m|;
\frac{r_0^2}{1+r_0^2}
\right)=0 .
\end{align}
For each fixed \((n,m)\), the corresponding branch satisfies
\begin{align}
\omega_{n,m}(r_0)
=\Delta+2n+|m|+\mathcal{O}(r_0^{2-2\Delta}),
\qquad r_0\to\infty .
\end{align}
Thus the finite cavity approaches the standard-quantization spectrum
\begin{align}
\omega_{n,m}^{(\mathrm{std})}=\Delta+2n+|m|.
\end{align}

\subsection{Robin Wall and Alternative Quantization}

For alternative quantization write \(\Delta=\Delta_-\) with
\begin{align}
0<\Delta\le1,\qquad M^2=\Delta(\Delta-2).
\end{align}
Adding the wall term
\begin{align}
S_{\Gamma_{r_0}}=-\frac{\Delta}{2}\int_{\Gamma_{r_0}}\dd^2x
\sqrt{-\gamma}\,\phi^2
\end{align}
gives the finite-radius Robin condition
\begin{align}
\left.(n^\mu\nabla_\mu+\Delta)\phi\right|_{r_0}=0,
\qquad n^r=\sqrt{1+r_0^2}.
\end{align}
For the regular radial branch this is
\begin{align}
Q_m^{\mathrm R}(\omega;r_0)
:=\left.\left(\sqrt{1+r^2}\,\pd_r+\Delta\right)
R_{\omega m}^{(\Delta)}(r)\right|_{r=r_0}=0 .
\end{align}
For \(0<\Delta<1\), the finite-radius branches obey
\begin{align}
\omega_{n,m}(r_0)
=\Delta+2n+|m|+\mathcal{O}(r_0^{-2\Delta}),
\qquad r_0\to\infty ,
\end{align}
and converge mode by mode to the alternative-quantization spectrum
\begin{align}
\omega_{n,m}^{(\mathrm{alt})}=\Delta+2n+|m|.
\end{align}
At the BF point \(\Delta=1\), the same Robin regulator selects the no-log
branch and gives
\begin{align}
\omega_{n,m}^{(\mathrm{BF})}=1+2n+|m|.
\end{align}

\section{Maxwell Field}

In temporal gauge the Gauss-law equation gives
\begin{align}
\pd_r f^r+\frac{1-r^2}{r(1+r^2)}f^r+\ii m f^\phi=0
\end{align}
for a mode \(A^\mu=e^{-\ii\omega t+\ii m\phi}f^\mu(r)\).  With
\[
f_r=\frac{f^r}{1+r^2},\qquad f_\phi=r^2 f^\phi ,
\]
the finite-wall boundary condition \(f^\phi(r_0)=0\) implies
\begin{align}
\pd_r(r f_r)\big|_{r_0}=0 .
\end{align}
Equivalently, write the physical temporal-gauge modes in terms of a scalar
potential \(\Psi\),
\begin{align}
A^t=0,\qquad
A^r=\frac{1+r^2}{r}\pd_\phi\Psi,\qquad
A^\phi=-\frac{1+r^2}{r}\pd_r\Psi .
\end{align}
Then \(A^\phi(r_0)=0\) is simply
\begin{align}
\pd_r\Psi\big|_{r_0}=0 .
\end{align}
For
\(\Psi=e^{-\ii\omega t+\ii m\phi}R_{\omega m}(r)\), the regular radial
branch is
\begin{align}
R_{\omega m}(r)=
r^{|m|}(1+r^2)^{-(|m|+2)/2}
{}_2F_1\!\left(
\frac{|m|+2-\omega}{2},
\frac{|m|+2+\omega}{2};
1+|m|;
\frac{r^2}{1+r^2}
\right).
\end{align}
The finite-radius Maxwell frequencies are therefore the positive roots of
\begin{align}
\pd_rR_{\omega m}(r)\big|_{r_0}=0 .
\end{align}
For \(m\ne0\), this is the same as the condition
\(\pd_r(r f_r)|_{r_0}=0\), since \(r f_r=\ii m\Psi\).  For \(m=0\), \(A^r=0\)
and \(A^\phi(r_0)=0\) again gives \(\pd_r\Psi|_{r_0}=0\).  For each fixed
\((n,m)\), the roots approach the infinite-radius temporal-gauge Maxwell
spectrum
\begin{align}
\omega_{n,m}^{(\infty)}=2+2n+|m| .
\end{align}

\section{Massive Proca Field}

The Proca action is
\begin{align}
S=\int\dd^3x\sqrt{-g}\left(
-\frac14 F_{\mu\nu}F^{\mu\nu}-\frac12\mu^2A_\mu A^\mu
\right),\qquad \mu^2>0 .
\end{align}
The equation of motion is
\begin{align}
\nabla_\nu F^{\nu\mu}-\mu^2 A^\mu=0 .
\end{align}
Taking the divergence gives the kinematic Proca constraint
\begin{align}
\nabla_\mu A^\mu=0 .
\end{align}
Since \(\sqrt{-g}=r\),
\begin{align}
\nabla_\mu A^\mu=
\pd_tA^t+\frac1r\pd_r(rA^r)+\pd_\phi A^\phi .
\end{align}
If the finite wall fixes \(A^t=A^\phi=0\) as functions of \((t,\phi)\), then
their tangential derivatives vanish on the wall.  The Proca constraint
therefore imposes the radial Robin condition
\begin{align}
\boxed{\;\pd_r(rA^r)\big|_{r_0}=0\;}
\end{align}
on every on-shell mode.  In terms of the covariant radial component
\(A_r=A^r/(1+r^2)\), this is
\begin{align}
\pd_r\!\left(r(1+r^2)A_r\right)\big|_{r_0}=0 .
\end{align}

\subsection{Decoupled Bulk Equations}

After eliminating \(A^t\) with the Proca constraint, define
\begin{align}
\psi=A^r,\qquad \varphi=rA^\phi,\qquad
\Phi_\pm=\frac12(\varphi\pm \ii\psi).
\end{align}
Then the two helicity variables satisfy
\begin{align}
\left[
(1+r^2)\pd_r^2+\frac{1+3r^2}{r}\pd_r
-\frac{1}{1+r^2}\pd_t^2
+1-\mu^2
+\frac{1}{r^2}\left(\pd_\phi^2\mp2\ii\pd_\phi-1\right)
\right]\Phi_\pm=0 .
\end{align}
For a mode \(e^{-\ii\omega t+\ii m\phi}\), the radial equations are scalar
equations with effective angular momenta
\begin{align}
\ell_+=|m-1|,\qquad \ell_-=|m+1|,
\end{align}
and effective scalar dimension
\begin{align}
\Delta=1+\mu .
\end{align}
The regular radial branch is
\begin{align}
R_\ell(\omega;r)=
r^\ell(1+r^2)^{-(\ell+\Delta)/2}
{}_2F_1\!\left(
\frac{\ell+\Delta-\omega}{2},
\frac{\ell+\Delta+\omega}{2};
1+\ell;
\frac{r^2}{1+r^2}
\right).
\end{align}

\subsection{Finite-Wall Spectrum}

For each fixed \(m\), the general regular solution is
\begin{align}
\Phi_+(r)=c_+R_{\ell_+}(\omega;r),\qquad
\Phi_-(r)=c_-R_{\ell_-}(\omega;r).
\end{align}
The wall conditions are
\begin{align}
\varphi(r_0)=0,\qquad \pd_r(r\psi)\big|_{r_0}=0 .
\end{align}
Equivalently,
\begin{align}
c_+R_{\ell_+}(r_0)+c_-R_{\ell_-}(r_0)&=0,\\
c_+\pd_r\!\left(rR_{\ell_+}\right)\big|_{r_0}
-c_-\pd_r\!\left(rR_{\ell_-}\right)\big|_{r_0}&=0 .
\end{align}
It is useful to package the finite-wall condition as the spectral function
\begin{align}
D_{r_0}(\omega)
&=R_{\ell_+}(r_0)\,\pd_r\!\left(rR_{\ell_-}\right)\big|_{r_0}
+R_{\ell_-}(r_0)\,\pd_r\!\left(rR_{\ell_+}\right)\big|_{r_0}.
\end{align}
Non-trivial finite-radius Proca modes are precisely the zeros of
\(D_{r_0}\).  The two terms in \(D_{r_0}\) show that the wall mixes the two
helicity variables.  A single helicity branch is not enough to satisfy both
\(A^\phi|_{r_0}=0\) and \(\pd_r(rA^r)|_{r_0}=0\), unless the complementary
boundary equation becomes redundant at a special frequency.

The infinite-volume limit is obtained from the large-\(r\) continuation of the
same regular branch.  For generic non-integer \(\mu\),
\begin{align}
R_\ell(\omega;r)
&=S_\ell(\omega)r^{\mu-1}+N_\ell(\omega)r^{-\mu-1}+\cdots,\\
S_\ell(\omega)
&=\frac{\Gamma(1+\ell)\Gamma(\mu)}
{\Gamma\!\left(\frac{\ell+\Delta-\omega}{2}\right)
\Gamma\!\left(\frac{\ell+\Delta+\omega}{2}\right)} .
\end{align}
Here \(S_\ell\) multiplies the non-normalizable asymptotic branch.  Substituting
this expansion into \(D_{r_0}\) gives the normalized spectral function
\begin{align}
F_{r_0}(\omega)
&\equiv r_0^{2-2\mu}D_{r_0}(\omega)\\
&=2\mu S_{\ell_+}(\omega)S_{\ell_-}(\omega)+o(1),
\end{align}
so that
\begin{align}
F_\infty(\omega)=2\mu S_{\ell_+}(\omega)S_{\ell_-}(\omega).
\end{align}
The large-\(r\) expansion, or its logarithmic limiting form for integer
\(\mu\), is uniform for \(\omega\) in every compact subset of the complex
plane.  Hence \(F_{r_0}\to F_\infty\) uniformly on compact \(\omega\)-sets.

For positive frequency, the zeros of \(F_\infty\) occur when one of the
non-normalizable coefficients vanishes,
\begin{align}
\frac{\ell_\pm+\Delta-\omega}{2}=-n,\qquad n=0,1,2,\ldots .
\end{align}
This gives the two infinite-volume Proca families
\begin{align}
\omega_{n,m,+}^{(\infty)}=\mu+1+2n+|m-1|,\qquad
\omega_{n,m,-}^{(\infty)}=\mu+1+2n+|m+1|.
\end{align}
The compact convergence gives the precise spectral statement: if \(K\) is a
compact domain whose boundary contains no zero of \(F_\infty\), then for
sufficiently large \(r_0\) the number of zeros of \(D_{r_0}\) in \(K\) equals
the number of zeros of \(F_\infty\) in \(K\), counted with multiplicity.  Thus
the fixed low-lying roots, counted in any finite frequency window, converge as
a multiset to the infinite-volume Proca spectrum.  There is no uniform
statement that all roots over the whole positive real axis stay bounded as
\(r_0\to\infty\), since the infinite-volume spectrum itself is unbounded as
\(n\to\infty\).

For example, for \(\mu=2\) and \(m=1\), the root near the first \(+\)-branch
value \(\omega=3\) is
\begin{align}
\omega(r_0=10)&=3.0003843828,\\
\omega(r_0=20)&=3.0000247511,\\
\omega(r_0=40)&=3.0000015586 .
\end{align}

\section{Mode Expansion Without Solving the Constraint}

At the action level one can avoid solving the Proca constraint first.  The
unreduced component action can be written as
\begin{align}
S=\frac12\int\dd t\,\dd r\,\dd\phi\Bigg[
&r\left(\frac{\dot A^r}{f}+\pd_r(fA^t)\right)^2
+\frac{1}{fr}\left(r^2\dot A^\phi+f\pd_\phi A^t\right)^2 \notag\\
&-\frac{f}{r}\left(\pd_r(r^2A^\phi)-\frac1f\pd_\phi A^r\right)^2
+\mu^2rf(A^t)^2-\frac{\mu^2r}{f}(A^r)^2-\mu^2r^3(A^\phi)^2
\Bigg],
\end{align}
where \(f=1+r^2\).  Choose componentwise regular spatial bases on
\([0,r_0]\), with \(A^t\) and \(A^\phi\) basis functions satisfying
Dirichlet boundary conditions at \(r_0\).  Expand at each time as
\begin{align}
A^t(t,r,\phi)&=\sum_I a_I(t)T_I(r,\phi),&
A^r(t,r,\phi)&=\sum_I b_I(t)R_I(r,\phi),&
A^\phi(t,r,\phi)&=\sum_I c_I(t)P_I(r,\phi).
\end{align}
Substitution gives a mechanical system
\begin{align}
S=\frac12\int\dd t\,
\left(K^{AB}_{IJ}\dot q_{I,A}\dot q_{J,B}
+2B^{AB}_{IJ}q_{I,A}\dot q_{J,B}
-V^{AB}_{IJ}q_{I,A}q_{J,B}\right),
\end{align}
with \(q_{I,A}=(a_I,b_I,c_I)\).  The \(a_I\) variables have no independent
kinetic term; their equations are the discretized Proca constraint.  Solving
or integrating out these auxiliary variables after the spatial expansion is
therefore equivalent to imposing the constraint after mode expansion.

The useful caution is that the spatial basis must be regular component by
component.  Splitting an on-shell highest-weight vector mode into independent
\(A^t,A^r,A^\phi\) coefficients is not a good off-shell basis: the separate
\(A^r\) and \(A^\phi\) pieces can be singular at the origin even though the full
on-shell vector is regular.  A finite-cavity off-shell calculation should
therefore use a componentwise regular Sturm--Liouville basis, not the separated
components of a highest-weight solution as independent off-shell variables.

\end{document}
